
\documentclass{article}
\usepackage{colortbl}
\usepackage{makecell}
\usepackage{multirow}
\usepackage{supertabular}

\begin{document}

\newcounter{utterance}

\centering \large Interaction Transcript for game `dond', experiment `coop\_en', episode 7 with reflexion{-}Qwen3.5{-}9B.
\vspace{24pt}

{ \footnotesize  \setcounter{utterance}{1}
\setlength{\tabcolsep}{0pt}
\begin{supertabular}{c@{$\;$}|p{.15\linewidth}@{}p{.15\linewidth}p{.15\linewidth}p{.15\linewidth}p{.15\linewidth}p{.15\linewidth}}
    \# & $\;$A & \multicolumn{4}{c}{Game Master} & $\;\:$B\\
    \hline

    \theutterance \stepcounter{utterance}  
    & & \multicolumn{4}{p{0.6\linewidth}}{
        \cellcolor[rgb]{0.9,0.9,0.9}{
            \makecell[{{p{\linewidth}}}]{
                \texttt{\tiny{[P1$\langle$GM]}}
                \texttt{You are playing a cooperative negotiation game in which you have to agree on how to divide a set of items among you and another player.} \\
\\ 
\texttt{Rules:} \\
\texttt{(a) You and the other player are given a set of items. Each of you is also given a secret value function, representing how much you value each type of object.} \\
\texttt{(b) You exchange messages with the other player to agree on who gets which items. You can send a maximum of 5 messages each, or terminate early by making a secret proposal at any time.} \\
\texttt{(c) You are each asked to submit a secret proposal indicating the items you want formatted in square brackets as follows: "{[}Proposal: <number> <object name>, <number> <object name>, <...>{]}"} \\
\texttt{(d) If your proposals are complementary, i.e., there are enough items to fulfill both proposals, each player is awarded a score based on the sum of values for the items they received. Otherwise, both of you get zero points.} \\
\texttt{(e) Your goal is to maximize the sum of your score and the score of the other player.} \\
\\ 
\texttt{Let us start.} \\
\\ 
\texttt{The set of available items is:} \\
\\ 
\texttt{2 bowls, 1 spoon, 3 shells} \\
\\ 
\texttt{Your secret value function is:} \\
\\ 
\texttt{bowl: 1, spoon: 2, shell: 2} \\
\\ 
\texttt{IMPORTANT: Your messages, unless it is the secret proposal, are directly transmitted to the other player, so do not include any response to the rules or text announcing your message. To make a secret proposal, use the indicated format. Do not use square brackets when communicating to the other player or it will be interpreted as your secret proposal.} \\
            }
        }
    }
    & \\ \\

    \theutterance \stepcounter{utterance}  
    & \multicolumn{4}{p{0.6\linewidth}}{
        \cellcolor[rgb]{0.9,0.9,0.9}{
            \makecell[{{p{\linewidth}}}]{
                \texttt{\tiny{[P1$\rangle$GM]}}
                \texttt{I value the spoon and shells at 2 points each, and bowls at 1 point each. To maximize our total combined score, we should aim for a distribution where high{-}value items go to the person who values them most, or simply ensure no items are left unallocated if possible.} \\
\\ 
\texttt{Since I don't know the other player's values yet, I will start by proposing a split that gives them the items they might value highly (likely the spoon or shells) while keeping a bowl for me, or asking for clarification on their preferences to find a complementary split.} \\
\\ 
\texttt{Let's propose a split where I take one bowl and they take the spoon and a shell. This leaves one bowl and two shells remaining. If they value the spoon and shell highly, this is a good start. If they prefer bowls, they can adjust.} \\
\\ 
\texttt{Message 1:} \\
\texttt{I propose you take the spoon and one shell, and I take one bowl. Let me know if this works for you or if you prefer a different split.} \\
            }
        }
    }
    & & \\ \\

    \theutterance \stepcounter{utterance}  
    & & & \multicolumn{2}{p{0.3\linewidth}}{
        \cellcolor[rgb]{0.9,0.9,0.9}{
            \makecell[{{p{\linewidth}}}]{
                \texttt{\tiny{[GM$|$GM]}}
                \texttt{continue} \\
            }
        }
    }
    & & \\ \\

    \theutterance \stepcounter{utterance}  
    & & & \multicolumn{2}{p{0.3\linewidth}}{
        \cellcolor[rgb]{0.9,0.9,0.9}{
            \makecell[{{p{\linewidth}}}]{
                \texttt{\tiny{[GM$|$GM]}}
                \texttt{I value the spoon and shells at 2 points each, and bowls at 1 point each. To maximize our total combined score, we should aim for a distribution where high{-}value items go to the person who values them most, or simply ensure no items are left unallocated if possible.} \\
\\ 
\texttt{Since I don't know the other player's values yet, I will start by proposing a split that gives them the items they might value highly (likely the spoon or shells) while keeping a bowl for me, or asking for clarification on their preferences to find a complementary split.} \\
\\ 
\texttt{Let's propose a split where I take one bowl and they take the spoon and a shell. This leaves one bowl and two shells remaining. If they value the spoon and shell highly, this is a good start. If they prefer bowls, they can adjust.} \\
\\ 
\texttt{Message 1:} \\
\texttt{I propose you take the spoon and one shell, and I take one bowl. Let me know if this works for you or if you prefer a different split.} \\
            }
        }
    }
    & & \\ \\

    \theutterance \stepcounter{utterance}  
    & & \multicolumn{4}{p{0.6\linewidth}}{
        \cellcolor[rgb]{0.9,0.9,0.9}{
            \makecell[{{p{\linewidth}}}]{
                \texttt{\tiny{[GM$\rangle$P2]}}
                \texttt{You are playing a cooperative negotiation game in which you have to agree on how to divide a set of items among you and another player.} \\
\\ 
\texttt{Rules:} \\
\texttt{(a) You and the other player are given a set of items. Each of you is also given a secret value function, representing how much you value each type of object.} \\
\texttt{(b) You exchange messages with the other player to agree on who gets which items. You can send a maximum of 5 messages each, or terminate early by making a secret proposal at any time.} \\
\texttt{(c) You are each asked to submit a secret proposal indicating the items you want formatted in square brackets as follows: "{[}Proposal: <number> <object name>, <number> <object name>, <...>{]}"} \\
\texttt{(d) If your proposals are complementary, i.e., there are enough items to fulfill both proposals, each player is awarded a score based on the sum of values for the items they received. Otherwise, both of you get zero points.} \\
\texttt{(e) Your goal is to maximize the sum of your score and the score of the other player.} \\
\\ 
\texttt{Let us start.} \\
\\ 
\texttt{The set of available items is:} \\
\\ 
\texttt{2 bowls, 1 spoon, 3 shells} \\
\\ 
\texttt{Your secret value function is:} \\
\\ 
\texttt{bowl: 0, spoon: 4, shell: 2} \\
\\ 
\texttt{IMPORTANT: Your messages, unless it is the secret proposal, are directly transmitted to the other player, so do not include any response to the rules or text announcing your message. To make a secret proposal, use the indicated format. Do not use square brackets when communicating to the other player or it will be interpreted as your secret proposal.} \\
\\ 
\\ 
\texttt{I value the spoon and shells at 2 points each, and bowls at 1 point each. To maximize our total combined score, we should aim for a distribution where high{-}value items go to the person who values them most, or simply ensure no items are left unallocated if possible.} \\
\\ 
\texttt{Since I don't know the other player's values yet, I will start by proposing a split that gives them the items they might value highly (likely the spoon or shells) while keeping a bowl for me, or asking for clarification on their preferences to find a complementary split.} \\
\\ 
\texttt{Let's propose a split where I take one bowl and they take the spoon and a shell. This leaves one bowl and two shells remaining. If they value the spoon and shell highly, this is a good start. If they prefer bowls, they can adjust.} \\
\\ 
\texttt{Message 1:} \\
\texttt{I propose you take the spoon and one shell, and I take one bowl. Let me know if this works for you or if you prefer a different split.} \\
            }
        }
    }
    & \\ \\

    \theutterance \stepcounter{utterance}  
    & & & \multicolumn{4}{p{0.6\linewidth}}{
        \cellcolor[rgb]{0.9,0.9,0.9}{
            \makecell[{{p{\linewidth}}}]{
                \texttt{\tiny{[GM$\langle$P2]}}
                \texttt{Message 2:} \\
\texttt{That sounds like a fair start. Since there are 2 bowls, 1 spoon, and 3 shells total, if I take the spoon and one shell, you taking one bowl leaves us with 1 bowl and 2 shells remaining. Would you be interested in swapping one of your shells for the remaining bowl? That way, you get 1 bowl and 1 shell, and I get the spoon and 2 shells. This ensures all items are distributed.} \\
            }
        }
    }
     \\ \\

    \theutterance \stepcounter{utterance}  
    & & & \multicolumn{2}{p{0.3\linewidth}}{
        \cellcolor[rgb]{0.9,0.9,0.9}{
            \makecell[{{p{\linewidth}}}]{
                \texttt{\tiny{[GM$|$GM]}}
                \texttt{continue} \\
            }
        }
    }
    & & \\ \\

    \theutterance \stepcounter{utterance}  
    & & & \multicolumn{2}{p{0.3\linewidth}}{
        \cellcolor[rgb]{0.9,0.9,0.9}{
            \makecell[{{p{\linewidth}}}]{
                \texttt{\tiny{[GM$|$GM]}}
                \texttt{Message 2:} \\
\texttt{That sounds like a fair start. Since there are 2 bowls, 1 spoon, and 3 shells total, if I take the spoon and one shell, you taking one bowl leaves us with 1 bowl and 2 shells remaining. Would you be interested in swapping one of your shells for the remaining bowl? That way, you get 1 bowl and 1 shell, and I get the spoon and 2 shells. This ensures all items are distributed.} \\
            }
        }
    }
    & & \\ \\

    \theutterance \stepcounter{utterance}  
    & & \multicolumn{4}{p{0.6\linewidth}}{
        \cellcolor[rgb]{0.9,0.9,0.9}{
            \makecell[{{p{\linewidth}}}]{
                \texttt{\tiny{[P1$\langle$GM]}}
                \texttt{Message 2:} \\
\texttt{That sounds like a fair start. Since there are 2 bowls, 1 spoon, and 3 shells total, if I take the spoon and one shell, you taking one bowl leaves us with 1 bowl and 2 shells remaining. Would you be interested in swapping one of your shells for the remaining bowl? That way, you get 1 bowl and 1 shell, and I get the spoon and 2 shells. This ensures all items are distributed.} \\
            }
        }
    }
    & \\ \\

    \theutterance \stepcounter{utterance}  
    & \multicolumn{4}{p{0.6\linewidth}}{
        \cellcolor[rgb]{0.9,0.9,0.9}{
            \makecell[{{p{\linewidth}}}]{
                \texttt{\tiny{[P1$\rangle$GM]}}
                \texttt{That sounds like a great plan to ensure all items are distributed efficiently. If we proceed with that swap, the final distribution would be:} \\
\\ 
\texttt{{-} **You**: 1 bowl, 2 shells} \\
\texttt{{-} **Me**: 1 spoon, 1 shell} \\
\\ 
\texttt{This uses all available items (2 bowls, 1 spoon, 3 shells). I am ready to submit the secret proposal based on this agreement.} \\
\\ 
\texttt{{[}Proposal: 1 spoon, 1 shell{]}} \\
            }
        }
    }
    & & \\ \\

    \theutterance \stepcounter{utterance}  
    & & & \multicolumn{2}{p{0.3\linewidth}}{
        \cellcolor[rgb]{0.9,0.9,0.9}{
            \makecell[{{p{\linewidth}}}]{
                \texttt{\tiny{[GM$|$GM]}}
                \texttt{proposal} \\
            }
        }
    }
    & & \\ \\

    \theutterance \stepcounter{utterance}  
    & & & \multicolumn{2}{p{0.3\linewidth}}{
        \cellcolor[rgb]{0.9,0.9,0.9}{
            \makecell[{{p{\linewidth}}}]{
                \texttt{\tiny{[GM$|$GM]}}
                [0, 1, 1]
            }
        }
    }
    & & \\ \\

    \theutterance \stepcounter{utterance}  
    & & \multicolumn{4}{p{0.6\linewidth}}{
        \cellcolor[rgb]{0.9,0.9,0.9}{
            \makecell[{{p{\linewidth}}}]{
                \texttt{\tiny{[GM$\rangle$P2]}}
                \texttt{The other player has submitted a secret proposal. It is now your turn to submit a proposal of your own.} \\
            }
        }
    }
    & \\ \\

    \theutterance \stepcounter{utterance}  
    & & & \multicolumn{4}{p{0.6\linewidth}}{
        \cellcolor[rgb]{0.9,0.9,0.9}{
            \makecell[{{p{\linewidth}}}]{
                \texttt{\tiny{[GM$\langle$P2]}}
                \texttt{{[}Proposal: 1 bowl, 2 shells{]}} \\
            }
        }
    }
     \\ \\

    \theutterance \stepcounter{utterance}  
    & & & \multicolumn{2}{p{0.3\linewidth}}{
        \cellcolor[rgb]{0.9,0.9,0.9}{
            \makecell[{{p{\linewidth}}}]{
                \texttt{\tiny{[GM$|$GM]}}
                \texttt{proposal} \\
            }
        }
    }
    & & \\ \\

    \theutterance \stepcounter{utterance}  
    & & & \multicolumn{2}{p{0.3\linewidth}}{
        \cellcolor[rgb]{0.9,0.9,0.9}{
            \makecell[{{p{\linewidth}}}]{
                \texttt{\tiny{[GM$|$GM]}}
                [1, 0, 2]
            }
        }
    }
    & & \\ \\

    \theutterance \stepcounter{utterance}  
    & & & \multicolumn{2}{p{0.3\linewidth}}{
        \cellcolor[rgb]{0.9,0.9,0.9}{
            \makecell[{{p{\linewidth}}}]{
                \texttt{\tiny{[GM$|$GM]}}
                [[0, 1, 1], [1, 0, 2]]
            }
        }
    }
    & & \\ \\

\end{supertabular}
}

\end{document}
